\section{انقباض گوسی}
فرض‌کنید که توابع 
$\phi_j : \mathbb{R} \to \mathbb{R}$
به ازای 
$j= 1, 2, \hdots, d$
،
\lr{1-Lipschitz}
باشند.
در این صورت نشان دهید که اگر تابع 
$\phi$
را به صورت روبرو تعریف کنیم:
$\phi(\mathbb{T}) := \{(\phi_1(\theta_1), \phi_1(\theta_1), \hdots, \phi_n(\theta_n) | \theta \in \mathbb{T} \} \subset \mathbb{R}^d$
که در آن 
$\mathbb{T} \in \mathbb{R}^d$
می‌باشد، نشان دهید که رابطه زیر بین پیچیدگی های گوسی برقرار خواهد بود:
$$\mathcal{G}(\phi(\mathbb{T})) \leq \mathcal{G}(\mathbb{T})$$
راهنمایی:
می‌توانید از قضیه \lr{Sudakov–Fernique}
استفاده کنید. همچنین فرض کنید که ابعاد بردار ها در این قضیه می‌تواند ناشمارا باشد.